% Options for packages loaded elsewhere
\PassOptionsToPackage{unicode}{hyperref}
\PassOptionsToPackage{hyphens}{url}
\documentclass[
]{article}
\usepackage{xcolor}
\usepackage{amsmath,amssymb}
\setcounter{secnumdepth}{-\maxdimen} % remove section numbering
\usepackage{iftex}
\ifPDFTeX
  \usepackage[T1]{fontenc}
  \usepackage[utf8]{inputenc}
  \usepackage{textcomp} % provide euro and other symbols
\else % if luatex or xetex
  \usepackage{unicode-math} % this also loads fontspec
  \defaultfontfeatures{Scale=MatchLowercase}
  \defaultfontfeatures[\rmfamily]{Ligatures=TeX,Scale=1}
\fi
\usepackage{lmodern}
\ifPDFTeX\else
  % xetex/luatex font selection
\fi
% Use upquote if available, for straight quotes in verbatim environments
\IfFileExists{upquote.sty}{\usepackage{upquote}}{}
\IfFileExists{microtype.sty}{% use microtype if available
  \usepackage[]{microtype}
  \UseMicrotypeSet[protrusion]{basicmath} % disable protrusion for tt fonts
}{}
\makeatletter
\@ifundefined{KOMAClassName}{% if non-KOMA class
  \IfFileExists{parskip.sty}{%
    \usepackage{parskip}
  }{% else
    \setlength{\parindent}{0pt}
    \setlength{\parskip}{6pt plus 2pt minus 1pt}}
}{% if KOMA class
  \KOMAoptions{parskip=half}}
\makeatother
\setlength{\emergencystretch}{3em} % prevent overfull lines
\providecommand{\tightlist}{%
  \setlength{\itemsep}{0pt}\setlength{\parskip}{0pt}}
\usepackage{bookmark}
\IfFileExists{xurl.sty}{\usepackage{xurl}}{} % add URL line breaks if available
\urlstyle{same}
\hypersetup{
  hidelinks,
  pdfcreator={LaTeX via pandoc}}

\author{}
\date{}

\usepackage[numbers]{natbib}
\begin{document}

\section{Understanding Vision-Language Models and Geometric Annotations
in X-ray
Imaging}\label{understanding-vision-language-models-and-geometric-annotations-in-x-ray-imaging}

\subsection{Formalized Questions}\label{formalized-questions}

\begin{itemize}
\tightlist
\item
  \textbf{Q1. Advances in Vision/Lang Models for X-ray Understanding:}
  What are the latest developments in large vision or vision-language
  models for interpreting X-ray images? (Including examples like
  \textbf{MedGemma} and similar models, their capabilities and
  applications in radiology.)
\item
  \textbf{Q2. Classical Annotation \& Measurements:} How are medical
  images (especially X-rays) traditionally annotated and measured
  geometrically? What ``geometric language'' or methods do clinicians
  use to mark images (e.g. measuring sizes, angles), and are these
  typically done manually? Are there automated methods for obtaining
  such measurements?
\item
  \textbf{Q3. Modern Approaches for Geometric Annotation (CAD-like):} Is
  there a modern approach that combines (1) and (2) -- for instance,
  using advanced AI models to produce formal geometric annotations of
  medical images (akin to a CAD model with shapes, splines, quantitative
  measurements)? In other words, do any large or foundational models
  exist that can perform detailed geometric \textbf{metrology} (shape
  measurements) on 2D X-ray images (e.g. in orthopedic, chest, or dental
  applications)?
\end{itemize}

With these clarified questions, we proceed to research the current state
of the art and literature.

\subsection{Advances in Vision and Vision-Language Models for X-ray
Image
Understanding}\label{advances-in-vision-and-vision-language-models-for-x-ray-image-understanding}

Large \textbf{vision} and \textbf{vision-language models (VLMs)} are
increasingly applied to medical imaging tasks, including X-ray
interpretation. These models leverage massive image and text data to
learn generalizable representations, enabling tasks like \textbf{image
captioning, report generation, visual question answering (VQA)}, and
classification in the medical
domain\cite{key-1}\cite{key-2}.
Notable recent advances include specialized models and open frameworks:

\begin{itemize}
\item
  \textbf{MedGemma (Google, 2025):} A state-of-the-art open multimodal
  model optimized for healthcare. MedGemma accepts medical images (e.g.
  chest X-rays) and text, and generates text (like radiology reports or
  answers to
  questions)\cite{key-3}.
  It comes in 4B and 27B parameter versions. Impressively,
  \textbf{MedGemma 4B} can generate chest X-ray reports that were judged
  sufficiently accurate for patient management in 81\% of cases by a
  board-certified
  radiologist\cite{key-4}.
  It also achieved competitive performance on medical image
  classification benchmarks, approaching task-specific
  models\cite{key-5}.
  MedGemma is part of Google's Health AI Developer Foundations (HAI-DEF)
  initiative, which emphasizes lightweight, privacy-preserving
  models\cite{key-6}\cite{key-7}.
  An associated image encoder called \textbf{MedSigLIP} (400M
  parameters) was trained on diverse medical images to power MedGemma's
  vision
  understanding\cite{key-8}\cite{key-9}.
  MedGemma has demonstrated strong results in chest X-ray report
  generation (RadGraph F1 of 30.3 after
  fine-tuning)\cite{key-10}
  and maintains general language capabilities, enabling it to handle
  mixed medical/non-medical queries.
\item
  \textbf{MedDR (2024):} A large medical vision-language model
  introduced by He et al. (2024) that uses \emph{diagnosis-guided
  bootstrapping} for
  training\cite{key-11}.
  MedDR was developed to improve disease detection across multiple
  imaging modalities. It represents a specialized VLM tailored for
  radiology diagnostics.
\item
  \textbf{LLaVA-Med (2024):} An adaptation of the LLaVA (Large
  Language-and-Vision Assistant) framework for
  medicine\cite{key-11}.
  LLaVA-Med is trained on medical image--text pairs, enabling a
  conversational AI to answer questions or explain findings from images.
  It leverages the success of multimodal conversational models (like
  GPT-4 Vision) but with domain-specific tuning.
\item
  \textbf{BiomedCLIP (Microsoft, 2023):} A medical version of the CLIP
  model, aligning radiology images with text reports or
  labels\cite{key-12}.
  BiomedCLIP provides a powerful image-text embedding; for example, it
  can classify diseases by computing similarities between an X-ray image
  embedding and text descriptions. Unlike generative VLMs, BiomedCLIP
  outputs probability scores for predefined
  categories\cite{key-13}\cite{key-14}.
  It has shown strong performance on some tasks (even surpassing larger
  VLMs on certain chest X-ray disease detection
  benchmarks)\cite{key-15}.
\item
  \textbf{Other Vision-Language Advances:} Research has also explored
  using general models like \textbf{GPT-4 with vision} for medical
  images. In fact, studies reported that GPT-4's vision system could
  achieve performance comparable to junior radiologists on certain image
  interpretation
  tasks\cite{key-16}.
  However, robustness remains an issue -- e.g. VLMs including GPT-4 and
  MedGemma can be sensitive to image artefacts (noise, motion blur,
  etc.), leading to drops in
  accuracy\cite{key-17}\cite{key-18}.
  This underscores ongoing challenges in reliability and the need for
  domain-specific fine-tuning for clinical
  use\cite{key-19}.
\end{itemize}

In summary, there is a surge of \textbf{foundation models} for medical
imaging. These include \textbf{multimodal generative models} (like
MedGemma) that produce free-text reports or answers, and
\textbf{image-text representation models} (like BiomedCLIP, MedSigLIP)
for classification and
retrieval\cite{key-9}.
They are transforming X-ray understanding by allowing AI to
\textbf{``see''} an X-ray and either describe it in natural language or
match it to clinical findings. The latest models focus on chest X-rays
(owing to large datasets like MIMIC-CXR), but efforts are broadening to
other radiographs (e.g. musculoskeletal X-rays for fractures or
alignment)\cite{key-20}\cite{key-21}.
Overall, these large models provide a foundation for interpreting
images, although their direct output is usually text or labels -- not
yet rich geometric annotations.

\subsection{Classical Annotation Methods and Geometric Measurements in
Medical
Imaging}\label{classical-annotation-methods-and-geometric-measurements-in-medical-imaging}

Traditionally, medical image annotations are largely \textbf{manual} and
rely on the expertise of clinicians (radiologists, orthopedists, etc.).
In daily practice, doctors use \textbf{PACS} software or dedicated tools
to mark images with regions of interest and perform measurements. Key
aspects of classical annotation include:

\begin{itemize}
\item
  \textbf{Annotation Types:} Common annotation techniques in medical
  images are \textbf{bounding boxes}, \textbf{polygons}, and
  \textbf{landmarks/keypoints} to outline structures or pinpoint
  anatomical
  landmarks\cite{key-22}.
  For example, a radiologist may draw a polygon around a tumor on an
  X-ray or mark key points on a bone's edges. These annotations serve as
  ground truth for AI training and aid in visualizing findings.
  Annotations often capture semantic labels (what structure or
  abnormality it is) and location/shape information.
\item
  \textbf{Geometric Measurements:} Beyond simple labeling, clinicians
  frequently take \textbf{quantitative measurements} from images. This
  includes linear dimensions (e.g. tumor diameter, bone length), angles
  (e.g. joint angles, spinal Cobb angle), areas, and ratios. Such
  \textbf{measurement annotations} are typically recorded as numeric
  values alongside the
  image\cite{key-23}.
  For instance, in a chest X-ray the cardiothoracic ratio (heart width
  to chest width) might be measured with digital calipers, or in an
  orthopedic X-ray the angle of a fracture or alignment of a joint is
  measured. These measurements are vital for diagnosis and treatment
  planning (e.g. monitoring tumor size or assessing scoliosis severity).
\item
  \textbf{Manual Process and Limitations:} The \textbf{gold standard}
  for image measurements is manual estimation by experts. For example,
  measuring the Cobb angle of a scoliotic spine on an X-ray involves a
  doctor identifying end-vertebrae and drawing lines -- a time-consuming
  process\cite{key-24}.
  Manual annotation is accurate when done carefully, but it is
  \textbf{labor-intensive and subject to inter-observer variability}.
  Studies note that manual Cobb angle measurement is arduous and can
  vary between observers, motivating automated
  solutions\cite{key-24}\cite{key-25}.
  Likewise, in dental cephalometric X-rays, orthodontists traditionally
  locate dozens of landmarks by hand to derive distances/angles, which
  is tedious and requires years of training.
\item
  \textbf{Tools and Languages:} There isn't a single formal ``geometric
  language'' universally used in annotations; instead, clinicians use
  imaging software features. \textbf{DICOM} (the standard for medical
  images) allows storing annotations and measurements (often as separate
  ``structured report'' objects or embedded presentation states). These
  might encode lines or spline contours drawn by the user. Essentially,
  the ``language'' of annotation is often geometric primitives: points,
  lines, circles, splines drawn on the image with associated labels or
  values. In classical practice, these annotations are not programmatic
  descriptions but visual marks and text notes on the image.
\end{itemize}

\textbf{Are these annotations automated?} Historically, \textbf{most
geometric annotation and measurement has been manual}, but recent years
have seen increasing automation:

\begin{itemize}
\item
  \textbf{Segmentation and Detection Algorithms:} Computer vision
  techniques have long attempted to automate anatomical measurements.
  Early methods (pre-deep-learning) included edge detection and
  model-based algorithms to find boundaries, but with limited general
  success. Today, \textbf{deep learning} is the primary driver of
  automation. For example, convolutional neural networks can be trained
  to \textbf{segment anatomical structures} (outline organs or lesions),
  from which measurements can be derived automatically (e.g. computing
  the area of a segmented lung nodule). Landmark-detection networks can
  identify key anatomical points to then calculate distances or angles.
\item
  \textbf{Automated Orthopedic Measurements:} A wealth of research shows
  that deep learning can replicate specific measurements. In spinal
  X-rays, numerous studies have developed models to automatically
  estimate the \textbf{Cobb angle} (the angle of spinal curvature) from
  just the X-ray
  image\cite{key-26}.
  A recent systematic review of such deep learning algorithms for Cobb
  angle found a pooled mean error around 2.99°, with segmentation-based
  methods performing better than landmark-based ones (≈2.4° vs 3.3°
  error)\cite{key-27}\cite{key-28}.
  This suggests AI can approach the precision of experts for scoliosis
  assessment. Other works achieved similar success: for instance, a deep
  learning model by Xu \emph{et al.} directly predicts Cobb angles from
  spine X-rays with performance comparable to
  radiologists\cite{key-29}.
\item
  \textbf{Examples of Automated Measurements:}
\item
  \emph{Lower Limb Alignment:} Deep networks have been developed to
  measure angles in knee/leg X-rays (e.g. hip-knee-ankle angle for leg
  alignment) fully
  automatically\cite{key-30}.
  One study achieved rapid and accurate detection of limb alignment on
  full-leg X-rays using deep
  learning\cite{key-30}.
\item
  \emph{Shoulder Angles:} Models can measure angles like the critical
  shoulder angle on radiographs, aiding orthopedic
  assessments\cite{key-31}.
\item
  \emph{Foot Anatomy:} A commercial AI (Milvue, France) was reported to
  \textbf{automate foot radiograph measurements}, such as podiatric
  angles, with high
  reliability\cite{key-32}.
  This deep learning solution significantly reduced the time needed for
  foot alignment evaluations compared to manual radiologist
  measurements\cite{key-32}.
\item
  \emph{Dental Cephalometrics:} AI systems now automatically detect
  cephalometric landmarks on lateral skull X-rays, allowing computation
  of standard orthodontic measurements (overbite, jaw angle, etc.).
  These have shown promising accuracy, nearing expert-level landmark
  localization\cite{key-33}
  (various studies in 2021--2023).
\item
  \emph{Chest X-ray Metrics:} Automation is also entering chest
  radiography. A notable example is \textbf{cardiothoracic ratio (CTR)}
  measurement. Qure.ai recently received FDA clearance for an AI tool
  (\textbf{qXR-CTR}) that automatically calculates the cardiothoracic
  ratio from frontal chest
  X-rays\cite{key-34}.
  This algorithm finds the borders of the heart and thorax on the X-ray
  and outputs the ratio, aiding rapid cardiomegaly (enlarged heart)
  screening\cite{key-35}\cite{key-36}.
  The automation improves consistency and may enable early detection of
  conditions like heart failure by quantifying heart size
  objectively\cite{key-37}.
\end{itemize}

In summary, while doctors \textbf{traditionally annotate and measure
X-rays manually}, drawing on their knowledge to interpret geometry,
there is a clear trend toward \textbf{automation}. Modern AI models --
often task-specific -- can now handle many measurements (lengths,
angles, areas) in X-rays with accuracy often on par with human experts.
These automated approaches usually involve a \textbf{deep learning model
that outputs either a mask (segmentation) or keypoints}, which are then
used to compute the desired metric. They greatly speed up the workflow
and reduce observer variability, though in critical cases clinicians
still verify the AI's suggestions. Notably, these automated measurement
tools are not yet \emph{universal}; each is usually trained for a
specific purpose (e.g. measuring one particular angle or detecting one
type of lesion).

\subsection{Toward CAD-like Geometric Modeling and Metrology with
AI}\label{toward-cad-like-geometric-modeling-and-metrology-with-ai}

The idea of a \textbf{formal geometric language or CAD-like model} for
medical images is intriguing. In engineering, CAD models describe
objects with precise geometric primitives (lines, arcs, splines) which
allow exact measurements. Here we ask: is there a modern approach
(perhaps leveraging foundation models) that can output such a structured
geometric representation for a medical image (like an X-ray), rather
than just labels or unstructured predictions?

\textbf{Current State -- Segmentation to Shapes:} Most deep learning
models for medical imaging output \emph{pixel-wise} results
(segmentation masks) or point estimates, not a parametric geometry.
However, researchers have started bridging this gap. One approach is to
constrain or convert segmentation outputs into \textbf{smooth,
parametric shapes}. For example, Williams \emph{et al.} (MICCAI 2021)
proposed an interactive segmentation framework where the CNN's
segmentation contour is represented as a \textbf{B-spline
curve}\cite{key-38}.
By using B-spline explicit active surfaces (a type of shape model),
their system ensures the outlined structure is smooth and anatomically
plausible, and it allows a human to adjust the spline in real-time if
needed\cite{key-38}.
This marriage of deep learning with \textbf{explicit geometry modeling}
is akin to a CAD-like representation -- the contour of an organ can be
expressed by control points and curves rather than a crude pixel mask.
Such parametric representations make it easier to compute precise
distances, radii, angles, etc., and to enforce prior knowledge about
shape continuity.

Another line of work integrates deep learning with \textbf{active shape
models} or statistical shape models. These methods build a deformable
model (often spline or mesh based) of a structure (say, a bone) and have
the AI predict the deformation parameters. The result is a vector
graphic outline of the target anatomy. This is still an active research
area and not yet widespread in clinical AI tools, but it demonstrates
the potential for CAD-like precision. By fitting models like splines or
ellipses to detected structures, one can directly obtain \textbf{metric
properties} (e.g. the major/minor axis lengths of an elliptical tumor
outline).

\textbf{Vision-Language Models with Structured Output:} Large VLMs like
MedGemma currently produce free-form text (e.g. a paragraph report)
rather than a formal geometric description. They excel at summarizing
findings and even mentioning measurements that they infer from the
image, but those measurements come as natural language (e.g. ``the tumor
is approximately 2 cm in diameter'') rather than as an interactive
model. As of now, there's \textbf{no known foundational model that
outputs a full CAD model or a formal geometry script from an X-ray}. The
focus has been on \textbf{semantic understanding} and language
generation, not on creating a quantifiable model of anatomy.

However, combining these advances is a logical next step. We can
envision a system where a foundation model interprets the image,
identifies the key structures, and then either: - calls a dedicated
module to fit geometric shapes to those structures, or - directly
generates a structured report containing coordinates or measurements.

Some early attempts exist in the form of \textbf{multi-step AI
pipelines}. For instance, an AI assistant could first use a segmentation
model (like a fine-tuned \textbf{Segment Anything Model}) to get a mask
of an organ or lesion, then compute geometric features from that mask,
and finally have a language model incorporate those features into its
report. This kind of pipeline isn't yet a one-click solution, but pieces
of it are in place. Indeed, the \textbf{Segment Anything Model (SAM)}, a
recent general model for segmentation, has been tested on medical images
(often with additional tuning) to quickly obtain masks of structures.
Those masks can yield volumes or areas which are the basis of metrology.
Researchers are actively exploring such combinations to enable
\textbf{quantitative analytics} in AI-driven diagnosis.

\textbf{Metrology in AI-Based Diagnosis:} ``Metrology'' refers to
quantitative measurement. In an AI diagnostic system, this could mean
the AI doesn't just say \emph{``there is a fracture''} but also
\emph{``the fracture is 3.5 cm in length with 10° angulation''}.
Achieving this requires both detection and measurement. Some current AI
products are moving in that direction on specific tasks (e.g., the
earlier example of Qure.ai's chest X-ray tool automatically calculating
a
ratio\cite{key-36}
-- here the output is a concrete number, not just a label). Similarly,
AI for fracture assessment might output the displacement in mm or the
fracture line length. These are early forms of integrating metrology
into AI results. They tend to rely on straightforward geometry derived
from segmentation or keypoints.

\textbf{Large/Foundational Models and Geometry:} We should note that
large pre-trained models can be fine-tuned or prompted to help with
measurements if given the right interfaces. For example, a
vision-language model could be prompted in a chain-of-thought manner to
\textbf{count pixels} corresponding to a known scale, effectively doing
a measurement calculation as part of its reasoning. But this is not
robustly reliable yet. No large medical VLM to date natively understands
a ``geometry language'' like a CAD program would. Instead, the precise
geometry understanding is usually achieved by dedicated vision models
(segmentation nets, etc.) as described.

In the \textbf{research frontier}, there are efforts to create more
\textbf{structured outputs}. One example is a foundation model called
\textbf{FLAnS (Free-form Language-based Segmentation)} which allows a
user to input text like ``segment the left lung and measure its area''
-- the model then segments and provides an output
mask\cite{key-39}.
Extending this concept, a future system might output a parametric
representation (for instance, return a polygon outline of the lung that
could be imported into a CAD software). Another example is integrating
\textbf{graph neural networks} or \textbf{geometric deep learning} to
handle anatomical structures (treating joints or vertebrae as nodes in a
graph, for instance) -- an approach that could inherently yield
geometric relations and thus
measurements\cite{key-40}\cite{key-41}.

\textbf{Summary of Modern Prospects:} We are beginning to see
convergence between \emph{vision-language understanding} and
\emph{geometric quantification}. While no single \textbf{X-ray
foundation model} currently outputs full CAD-like annotations,
researchers are introducing geometry-aware AI methods. By combining
segmentation networks, shape modeling (splines/meshes), and the
interpretive power of VLMs, one can imagine a ``\textbf{next-generation
CAD (Computer-Aided Diagnosis)}'' system in radiology that goes beyond
highlighting an abnormality -- it could \emph{draw} it out with precise
contours and numbers. For now, the most advanced systems remain
task-specific (focused on one type of measurement or one anatomy).
Foundation models like MedGemma provide the high-level understanding and
language
interface\cite{key-7},
whereas classical geometric methods (augmented with deep learning)
provide the numeric rigor. Merging these effectively is an open research
area.

\subsection{Conclusion}\label{conclusion}

In conclusion, \textbf{large vision and vision-language models} are
making significant strides in interpreting X-ray images. Models like
MedGemma exemplify how far we've come in generating accurate radiology
reports and answering clinical questions from
images\cite{key-4}.
At the same time, the routine work of \textbf{annotating and measuring
medical images} has historically been manual, using simple geometric
tools -- but AI is rapidly automating many of these measurements, from
spinal curvature
angles\cite{key-24}\cite{key-26}
to heart size
ratios\cite{key-34}.
The fusion of these two threads -- semantic understanding and precise
geometric quantification -- is on the horizon. Early research that
incorporates \textbf{formal geometric representations} (like
spline-based
contours\cite{key-38})
into deep learning shows promise for achieving CAD-like accuracy in
measurements.

\textbf{As of now, no single ``foundational'' AI system can intake an
X-ray and output a full geometric analysis (e.g. a complete CAD model of
a bone or tumor with all measurements)}. Instead, the capability is
split among specialized components: vision-language models excel at
description and reasoning, while dedicated vision models handle exact
localization and measurement. Modern AI in medical imaging is moving
toward uniting these capabilities. In practice, AI-based tools (some
FDA-approved) are already performing quantitative tasks that augment
diagnosis\cite{key-36}
-- essentially early steps of metrological analysis. It's reasonable to
expect that future foundational models or multi-model systems will
further integrate this, enabling clinicians not only to detect
abnormalities but also to obtain reliable \textbf{quantitative geometry}
(sizes, shapes, and other metrics) from X-ray images at the click of a
button.



\bibliographystyle{unsrtnat}
\bibliography{deep-research02}

\end{document}
